\documentclass[11pt]{amsart}
\usepackage{geometry}               
\geometry{letterpaper}                    

\usepackage{graphicx}
\usepackage{tikz}
\usepackage{amssymb}
\usepackage{pdfsync}
\usepackage{hyperref}
\usepackage{verbatim} % this enables the {comment} environment
\usepackage{epstopdf}
\DeclareGraphicsRule{.tif}{png}{.png}{`convert #1 `dirname #1`/`basename #1 .tif`.png}

\newcommand{\pvector}[1]{
  \begin{pmatrix}
    #1
  \end{pmatrix}} 
\newcommand{\ddirac}[1]{
  \,\boldsymbol{\delta}\!\pvector{#1}\!} 
\renewcommand{\d}{\,{\rm d}} 

\def\R{\mathbb{R}}
\def\T{\mathbb{T}}
\def\N{\mathbb{N}}
\def\Z{\mathbb{Z}}
\def\Co{\mathbb{C}}

\def\C{\mathcal{C}}
\def\Q{\mathcal{Q}}
\def\D{\mathcal{D}}
\def\E{\mathcal{E}}
\def\G{\mathcal{G}}
\def\P{\mathbb{P}}
\def\H{\mathcal{H}}
\def\F{\mathcal{F}}
\def\V{\mathcal{V}}
\def\X{\mathcal{X}}
\def\f{\widehat{f}}

\newtheorem{theorem}{Theorem}
\newtheorem{corollary}[theorem]{Corollary}
\newtheorem{definition}[theorem]{Definition}
\newtheorem{remark}[theorem]{Remark}
\newtheorem{proposition}[theorem]{Proposition}
\newtheorem{lemma}[theorem]{Lemma}
\newtheorem{estimate}[theorem]{Estimate}
\newtheorem{claim}[theorem]{Claim}
\newtheorem{question}[theorem]{Question}
\newtheorem{example}[theorem]{Example}
\newtheorem{conclusion}[theorem]{Conclusion}
\newtheorem*{notation}{Notation}


\title{An inverse problem for convolution 
measures with applications}




\author{Diogo Oliveira e Silva}
\address{
 Departamento de Matemática \\
Instituto Superior Técnico \\
Av. Rovisco Pais \\
1049-001 Lisboa, Portugal}
\email{diogo.oliveira.e.silva@tecnico.ulisboa.pt}

\author{João Pedro Ramos}
\address{
 EPFL Lausanne -- Institute of Mathematics MA C3 535 \\
Route Cantonale \\ 
1015 Lausanne, Switzerland}
\email{joao.ramos@epfl.ch}


\date{\today}      
\subjclass[2020]{42B10}
\keywords{Sharp restriction theory, Monge--Ampère equation, stability, inverse problem, convolution measure.} 

\begin{document}
\begin{abstract}
\end{abstract}

\maketitle

\section{Introduction}

A classical theorem of Jörgens ($d=2$), Calabi ($d\leq 5$) and Pogorelov
($d\geq 2$) states that any classical convex solution of
$\det D^2u=1$ in $\R^d$ must be a quadratic polynomial.  As stated in the
second paragraph of \cite{CL03}, Caffarelli extended the result from
classical solutions to {\it viscosity solutions}.  We recall the definitions.
\begin{definition}[\cite{Gu16}]
    Let $u\in C^0(\Omega)$ be a convex function and let
    $0\leq f\in C^0(\Omega)$.  The function $u$ is a viscosity subsolution
    (resp.\@ supersolution) of the equation $\det D^2 u=f$ in $\Omega$ if
    whenever convex $\phi\in C^2(\Omega)$ and $x_0\in\Omega$ are such that
    $(u-\phi)(x)\leq (u-\phi)(x_0)$ (resp.\@ $\geq$) for all $x$ in a
    neighborhood of $x_0$, then necessarily
    \[
        \det D^2\phi(x_0)\geq f(x_0)
        \qquad(\text{resp. } \leq).
    \]
\end{definition}
\noindent As in \cite[Remark 1.3.3]{Gu16}, no generality is lost in assuming
that $\phi(x_0)=u(x_0)$, and one may further restrict the test functions to
strictly convex quadratic polynomials.

Let $\psi:\R^d\to\R$ be convex, and let $\sigma$ denote projection
measure on the graph
\[
    \{(x,\psi(x)):x\in\R^d\}\subset\R^{d+1}.
\]
Thus
\[
    \langle \sigma,\varphi\rangle
    =\int_{\R^d}\varphi(x,\psi(x))\,\d x.
\]
The centered convolution quantity under scrutiny is
\[
    F_\sigma(\xi,\tau)
    :=(\sigma\ast\sigma)(\xi,2\psi(\xi/2)+\tau),\qquad \tau>0.
\]
Via the usual change of variables from \cite[Prop. 2.1(c)]{OSQ18} and
\cite[Prop. 2.1(c)]{OSQ20}, this can be written as
\begin{equation}\label{eq:centered-conv}
    F_\sigma(\xi,\tau)
    =\int_{\R^d}
        \ddirac{\tau-\Psi(\xi,y)}
      \,\d y,
    \qquad
    \Psi(\xi,y):=\psi(\xi/2+y)+\psi(\xi/2-y)-2\psi(\xi/2).
\end{equation}
Convexity gives $\Psi\geq 0$.  If $\psi$ is strictly convex, then
$\Psi(\xi,y)>0$ whenever $y\neq 0$.

For a quadratic polynomial
\[
    \psi(x)=a+b\cdot x+\frac12 x^T Hx,\qquad H=H^T>0,
\]
one has
\[
    \Psi(\xi,y)=y^THy
\]
for every $\xi,y$.  Consequently
\[
    F_\sigma(\xi,\tau)
    =\frac1{\sqrt{\det H}}\int_{\R^d}\ddirac{\tau-|w|^2}\,\d w
    =\frac{d\,\omega_d}{2\sqrt{\det H}}\tau^{d/2-1},
\]
where $\omega_d=|\{x\in\R^d:|x|<1\}|$.  Thus the constant profile from the
two-dimensional paraboloid is a special feature of dimension two; in higher
dimensions a quadratic graph has the power-law profile $\tau^{d/2-1}$.

The point of the argument below is that the integrated form of the convolution
identity is stronger than the almost everywhere Hessian information discussed
in the transcript.  It tests $\psi$ against upper and lower quadratic supports
and therefore gives a viscosity Monge--Ampère equation directly, once the
centered sections are known to shrink to their center.  This last point is
automatic under strict convexity, but we isolate it below because it is the only
place where strict convexity is used.

\section{An inverse result}

\begin{definition}\label{def:constant-normalized}
Let $c_0>0$.  We say that the centered convolution is constant on its
normalized support with constant $c_0$ if, for every $\xi\in\R^d$ and every
$t>0$,
\begin{equation}\label{eq:section-law}
    |\{y\in\R^d:\Psi(\xi,y)<t\}|
    =c_0t.
\end{equation}
Equivalently, after differentiating in $t$ in the sense of distributions for
each fixed $\xi$,
\begin{equation}\label{eq:profile}
    F_\sigma(\xi,\tau)=c_0,\qquad \tau>0.
\end{equation}
\end{definition}

\begin{definition}\label{def:shrinking}
We say that $\psi$ has shrinking centered sections if, for every $x\in\R^d$
and every $r>0$,
\[
    \min_{|y|=r}\bigl(\psi(x+y)+\psi(x-y)-2\psi(x)\bigr)>0.
\]
This is automatic if $\psi$ is strictly convex.
\end{definition}

\begin{theorem}[Constant normalized support]\label{thm:main}
Let $d\geq 1$, let $\psi:\R^d\to\R$ be convex, and let $\sigma$ be projection
measure on its graph.  Assume that the centered convolution is constant on its
normalized support with constant $c_0>0$.  Assume moreover that:
\begin{enumerate}
    \item $\psi$ has a nondegenerate Alexandrov second-order expansion at
    least one point;
    \item $\psi$ has shrinking centered sections.
\end{enumerate}
Then $d=2$, and
\[
    \psi(x)=a+b\cdot x+\frac12 x^T Hx
\]
for some $a\in\R$, $b\in\R^d$, and some positive definite symmetric matrix
$H$, with
\begin{equation}\label{eq:det-constant}
    \det H=\left(\frac{\pi}{c_0}\right)^2.
\end{equation}
After subtracting the affine part and making the linear change of variables
which sends $H$ to $2I$, this is exactly the model $\psi(x)=|x|^2$.

Conversely, every such positive definite quadratic polynomial in dimension two
satisfies \eqref{eq:profile} with $c_0=\pi/\sqrt{\det H}$.  In particular,
$\psi(x)=|x|^2$ gives $c_0=\pi/2$.
\end{theorem}

\begin{remark}
The exact conclusion $\psi(x)=|x|^2$ cannot follow from constancy alone unless
one fixes the affine and linear normalizations.  Indeed every positive
definite quadratic polynomial in dimension two has constant centered
convolution; the constant only sees $\det H$, not the eccentricity of $H$.
\end{remark}

\begin{lemma}[Local section asymptotics]\label{lem:local-section}
Let $\phi\in C^2$ be convex in a neighborhood of $x_0\in\R^d$, and set
\[
    Q(y):=\phi(x_0+y)+\phi(x_0-y)-2\phi(x_0),\qquad
    A:=D^2\phi(x_0).
\]
If $A$ is positive definite, then, for every $r>0$,
\begin{equation}\label{eq:local-asymp-pd}
    \lim_{t\to0^+}
    t^{-d/2}|\{y\in B(0,r):Q(y)<t\}|
    =\frac{\omega_d}{\sqrt{\det A}}.
\end{equation}
If $A$ is singular, then, for every $r>0$,
\begin{equation}\label{eq:local-asymp-sing}
    \lim_{t\to0^+}
    t^{-d/2}|\{y\in B(0,r):Q(y)<t\}|=+\infty
\end{equation}.
\end{lemma}

\begin{proof}
The positive definite case follows from
\[
    Q(y)=y^TAy+o(|y|^2)
\]
and the linear change of variables $w=A^{1/2}y$; for $t>0$ sufficiently small,
the relevant section is contained in $B(0,r)$.

If $A$ is singular, fix $r>0$ and split $\R^d=E\oplus K$, where
$K=\ker A$ has dimension $k\geq 1$ and $E=K^\perp$.  Since
$Q(y)=y^TAy+o(|y|^2)$, for every $\varepsilon>0$ there is a radius
$0<\rho<r$ such that, for $y=e+z\in E\oplus K$ with $|y|<\rho$,
\[
    Q(y)\leq 2\,e^TAe+\varepsilon |z|^2.
\]
For $t>0$ sufficiently small, the set
\[
    \{e\in E:2e^TAe<t/2\}\times
    \{z\in K:\varepsilon |z|^2<t/2\}
\]
is contained in $\{Q<t\}\cap B(0,r)$ and has measure
$C_A\varepsilon^{-k/2}t^{d/2}$.  Since $\varepsilon$ is arbitrary, the
normalized local section volume tends to $+\infty$.
\end{proof}

\section{Proof of the characterization}

\begin{proof}[Proof of Theorem~\ref{thm:main}]
We first pin down the dimension.  Let $u_0$ be a point at which $\psi$ has a
nondegenerate Alexandrov second-order expansion,
\[
    \psi(u_0+y)=\psi(u_0)+p\cdot y+\frac12 y^TH_0y+o(|y|^2),
    \qquad H_0=H_0^T>0.
\]
Then
\[
    \Psi(2u_0,y)=y^TH_0y+o(|y|^2),
\]
and therefore
\[
    |\{y\in\R^d:\Psi(2u_0,y)<t\}|
    \sim \frac{\omega_d}{\sqrt{\det H_0}}t^{d/2}
    \qquad (t\to0^+).
\]
Comparison with \eqref{eq:section-law} forces $d=2$ and
\[
    \det H_0=\left(\frac{\pi}{c_0}\right)^2.
\]

It remains to prove that $\psi$ is quadratic.  Set
\[
    \lambda:=\left(\frac{\pi}{c_0}\right)^2.
\]
We prove that $\psi$ is a viscosity solution of
\begin{equation}\label{eq:MA}
    \det D^2\psi=\lambda
    \qquad\text{on }\R^2.
\end{equation}

Let $\phi\in C^2$ be convex near $x_0$, with
\[
    \phi(x_0)=\psi(x_0),\qquad \phi\geq\psi
\]
near $x_0$.  Choose $r>0$ so small that the comparison holds whenever
$|y|<r$.  For such $y$,
\[
    \Psi(2x_0,y)
    \leq
    \phi(x_0+y)+\phi(x_0-y)-2\phi(x_0)=:Q(y).
\]
Hence, for every $t>0$,
\[
    \{y\in B(0,r):Q(y)<t\}
    \subset \{y\in\R^2:\Psi(2x_0,y)<t\}.
\]
Thus
\[
    |\{y\in B(0,r):Q(y)<t\}|\leq c_0t.
\]
In view of Lemma~\ref{lem:local-section}, the matrix $A=D^2\phi(x_0)$ cannot
be singular.  Therefore \eqref{eq:local-asymp-pd} gives
\[
    \frac{\pi}{\sqrt{\det A}}\leq c_0,
\]
or equivalently
\[
    \det D^2\phi(x_0)\geq \lambda.
\]
This is the viscosity subsolution inequality.

Now let $\phi\in C^2$ be convex near $x_0$, with
\[
    \phi(x_0)=\psi(x_0),\qquad \phi\leq\psi
\]
near $x_0$.  Choose $r>0$ so small that the comparison holds whenever
$|y|<r$.  By the shrinking-section assumption,
\[
    m_r:=\min_{|y|=r}\Psi(2x_0,y)>0,
\]
and $\{\Psi(2x_0,\cdot)<t\}\subset B(0,r)$ whenever $0<t<m_r$.  Indeed,
any $y$ with $|y|>r$ and $\Psi(2x_0,y)<m_r$ would give, by convexity along
the ray from $0$ to $y$,
\[
    \Psi(2x_0,r y/|y|)
    \leq \frac{r}{|y|}\Psi(2x_0,y)<m_r,
\]
which is impossible.  Hence
\[
    \{y\in\R^2:\Psi(2x_0,y)<t\}
    \subset \{y\in B(0,r):Q(y)<t\}
\]
whenever $0<t<m_r$.
If $A=D^2\phi(x_0)$ is singular, then $\det A=0\leq \lambda$.  If $A$ is
positive definite, then
\[
    c_0t\leq |\{y\in B(0,r):Q(y)<t\}|
\]
for $0<t<m_r$, and \eqref{eq:local-asymp-pd} implies
\[
    c_0\leq \frac{\pi}{\sqrt{\det A}},
\]
or equivalently
\[
    \det D^2\phi(x_0)\leq \lambda.
\]
This is the viscosity supersolution inequality.

Thus $\psi$ is a convex viscosity solution of \eqref{eq:MA}.  By the
Jörgens theorem in the viscosity form of Caffarelli--Li
\cite[Theorem 1.1]{CL03}, after the harmless normalization of the right-hand
side of \eqref{eq:MA} to $1$, $\psi$ is a quadratic polynomial.  This proves
the direct implication and the determinant identity \eqref{eq:det-constant}.

For the converse, let
\[
    \psi(x)=a+b\cdot x+\frac12 x^THx,\qquad H=H^T>0.
\]
Then $\Psi(\xi,y)=y^THy$, independently of $\xi$.  Consequently
\[
    F_\sigma(\xi,\tau)
    =\int_{\R^2}\ddirac{\tau-y^THy}\,\d y
    =\frac{\pi}{\sqrt{\det H}},
\]
which is \eqref{eq:profile} with
$c_0=\pi/\sqrt{\det H}$.
\end{proof}

\section{The degenerate alternative}

We now spell out what should replace the nondegeneracy hypothesis in higher
dimensions.  The dimension reduction in the proof of Theorem~\ref{thm:main}
used only one positive definite Alexandrov second-order expansion.  In
dimension $d\geq 3$ such a point is impossible under the linear section law
\eqref{eq:section-law}: if
\[
    \psi(x_0+y)=\psi(x_0)+p\cdot y+\frac12 y^TH_0y+o(|y|^2),
    \qquad H_0=H_0^T>0,
\]
then
\[
    |\{y\in\R^d:\Psi(2x_0,y)<t\}|
    \sim \frac{\omega_d}{\sqrt{\det H_0}}t^{d/2},
    \qquad t\to0^+,
\]
which is incompatible with $c_0t$ when $d\neq 2$.

In the rest of this section we also impose the lower normalization
\[
    \psi\geq0.
\]
This is the normalization relevant to the final geometric obstruction.

Thus, in dimensions $d\geq 3$, the only possible convex case is the completely
degenerate Alexandrov alternative:
\[
    \det D^2\psi(x)=0
    \quad\text{at every Alexandrov point }x.
\]
What is needed here is a developability statement for convex graphs.  In the
smooth complete case, the classical mechanism is the Hartman--Nirenberg
theorem: a complete $C^2$ hypersurface of constant zero Gauss--Kronecker
curvature in Euclidean space is cylindrical, i.e. it is ruled by affine
$(d-1)$-planes \cite{HN59}.  For nonsmooth convex hypersurfaces, the natural
language is Aleksandrov--Pogorelov bounded extrinsic curvature and the
associated spherical-image curvature measure; see Pogorelov's monograph
\cite{Po73}.  Ushakov gives a modern smooth account of developable surfaces in
arbitrary dimension and codimension \cite{Us99}.  The expected form needed
here is:
\begin{quote}
If a finite convex function $\psi:\R^d\to\R$ has vanishing extrinsic curvature
measure, and its graph is not an affine hyperplane, then through every point
$(x,\psi(x))$ of the graph there is a nontrivial affine segment; under the
global hypotheses relevant here, this segment should promote to a complete
affine line in the graph.
\end{quote}
The curvature-measure formulation is important.  The a.e. Alexandrov
degeneracy condition alone does not exclude singular curvature mass; it is
therefore too weak to force rulings through every point.
Equivalently, for every $x\in\R^d$ one should find a nonzero direction
$v_x\in\R^d$ such that the one-dimensional slice
\[
    s\mapsto \psi(x+sv_x)
\]
is affine on all of $\R$.  In that case
\[
    \Psi(2x,sv_x)=0\qquad\text{for every }s\in\R.
\]
This by itself is not yet a contradiction, since a line has zero
$d$-dimensional measure.  The required final step is a quantitative
``thickening of rulings'' argument: the presence of a complete flat direction
through each point of the graph should force the centered sections
\[
    \{y:\Psi(2x,y)<t\}
\]
to contain enough transverse thickness around that line to have either
infinite volume or a small-$t$ power law different from $t$.  In the ideal
cylindrical case, for example, the section would contain a product of the flat
line with a transverse $(d-1)$-dimensional section, so the volume would be
infinite.  In less rigid ruled situations, one should still aim to prove a
lower bound contradicting the exact linear identity \eqref{eq:section-law}.

This would finish the desired higher-dimensional contradiction.  More
precisely, the remaining task is to supply or cite the appropriate
developability theorem in the low regularity class at hand.  In the $C^2$
case this is the familiar flat-direction argument: $\det D^2\psi\equiv0$ and
convexity force a ruling by affine directions.  For merely convex functions,
one has to rule out singular Monge--Ampère mass, justify that Alexandrov
degeneracy gives actual affine rulings of the graph rather than only a.e. flat
directions, and then prove the quantitative thickening estimate for the
centered sections.

\subsection{The convolution identity forces the zero-curvature hypothesis}
There is a slightly better way to phrase the remaining low-regularity
assumption.  We should not assume that the curvature measure vanishes.  We
should assume only that $\psi$ belongs to a class in which the appropriate
curvature measure is defined, and then derive its vanishing from
\eqref{eq:section-law}.

Here is the argument.  Let $d\geq3$, fix $x_0\in\R^d$, and let
$\phi\in C^2$ be convex near $x_0$, with
\[
    \phi(x_0)=\psi(x_0),\qquad \phi\leq\psi
\]
near $x_0$.  As in the viscosity supersolution part of the proof of
Theorem~\ref{thm:main}, assume that the centered sections shrink to their
center, so that for $t>0$ sufficiently small the set
$\{\Psi(2x_0,\cdot)<t\}$ is contained in the neighborhood where the comparison
with $\phi$ holds.  Then
\[
    \{\Psi(2x_0,\cdot)<t\}
    \subset
    \{y:Q(y)<t\},
    \qquad
    Q(y):=\phi(x_0+y)+\phi(x_0-y)-2\phi(x_0).
\]
If $A=D^2\phi(x_0)$ were positive definite, Lemma~\ref{lem:local-section}
would give
\[
    |\{y:Q(y)<t\}|=O(t^{d/2}),
    \qquad t\to0^+.
\]
This contradicts \eqref{eq:section-law}, because $d/2>1$ and
$c_0t\not=O(t^{d/2})$.  Hence every lower quadratic test has singular Hessian:
\[
    \det D^2\phi(x_0)=0.
\]
In other words, $\psi$ is a convex viscosity supersolution of
\[
    \det D^2\psi=0.
\]
Since the Monge--Ampère measure of a convex function is nonnegative, this is
equivalent, in the Aleksandrov/viscosity correspondence, to vanishing
Monge--Ampère measure:
\[
    M_\psi=0.
\]
Thus, if we work in a class where the graph curvature measure is the
Aleksandrov--Pogorelov spherical-image measure, or in a class where
${\rm Det}\,D^2\psi$ is defined weakly, the vanishing condition is not an extra
hypothesis.  It is forced by the convolution identity, together with the
shrinking-section localization.  For graphs of convex functions, vanishing of
$M_\psi$ is the same zero-curvature condition after applying the Gauss-map
parametrization of the spherical image; the weight relating the two measures is
positive on bounded subgradient ranges.

\subsection{Possible hypotheses below \texorpdfstring{$C^1$}{C1}}
The preceding paragraph suggests the following hierarchy of assumptions.  In
the first two cases, the condition imposed at the outset is only membership in
the relevant weak class; the zero-curvature equation is then a consequence of
the convolution identity.  The last option is not below $C^1$ in its Hölder
version, and is best kept as a separate remark about a regime where one might
also be able to proceed.

\begin{enumerate}
    \item {\it Convex surface of bounded extrinsic curvature.}  Work in the
    Aleksandrov--Pogorelov category of convex hypersurfaces with bounded
    extrinsic curvature.  In graph language, this means that the relevant
    Gauss--Kronecker curvature measure of the graph is defined.  The preceding
    viscosity argument then forces this measure to vanish.  This is the closest
    low-regularity version of the Hartman--Nirenberg/Pogorelov mechanism; it
    allows corners a priori and is therefore genuinely below $C^1$.

    In this class, Pogorelov's developability theorem applies directly: a
    convex hypersurface of bounded extrinsic curvature whose curvature measure
    vanishes is developable.  In the graph setting this means the following.
    For every point $(x,\psi(x))$ of the graph there is a nontrivial affine
    segment contained in the graph and passing through $(x,\psi(x))$.  If the
    graph is complete and one uses the global form of the theorem, the local
    segments are the generators of a ruled hypersurface; away from the affine
    pieces these generators can be continued until they either meet the boundary
    or become complete affine lines.  Since our graph is entire, there is no
    finite boundary in the base, so this is the natural route to the desired
    complete rulings.

    \item {\it Monge--Ampère function with generalized determinant.}  Require
    that $\psi$ belongs to a class of Monge--Ampère functions in which the
    Hessian determinant is defined weakly and continuously.  The argument above
    should then give
    \[
        {\rm Det}\,D^2\psi=0
    \]
    as a measure/current, rather than requiring this as an assumption.  In
    dimension two, Jerrard's rigidity theorem for Monge--Ampère functions gives
    developability under such a weak formulation \cite{Je10}.  The point for us
    would be to use the analogous higher-dimensional determinant or
    Monge--Ampère system, if available, as a substitute for pointwise
    differentiability.

    Even without quoting the full Pogorelov language, one gets the local
    ruled conclusion from the Aleksandrov maximum principle.  Indeed, after
    subtracting a supporting affine function at $x_0$, suppose that the contact
    set of $\psi$ with that support were reduced to $x_0$ inside a small ball.
    Then a small section
    \[
        S_h:=\{x:\psi(x)<h\}
    \]
    would be a bounded convex set with $x_0\in S_h$ and $\psi=h$ on
    $\partial S_h$.  Aleksandrov's maximum principle bounds the height of such
    a section from above by a dimensional constant times the Monge--Ampère mass
    of $S_h$.  Since ${\rm Det}\,D^2\psi=0$, this mass is zero, forcing
    $h=0$, a contradiction.  Thus every contact set contains a second point,
    and the convexity of the contact set gives an affine segment in the graph
    through $(x_0,\psi(x_0))$.

    This proves the local ruling statement in the weak Monge--Ampère class.  To
    promote these local contact segments to complete affine lines in the full
    $d\geq3$ problem, one still needs the corresponding global developability
    theorem in that weak class.  Jerrard supplies this mechanism in dimension
    two; a higher-dimensional analogue would give the required replacement for
    Pogorelov in option (1).
\end{enumerate}

\begin{proposition}[Local rulings from zero Monge--Ampère measure]
\label{prop:local-rulings-ma}
Let $\Omega\subset\R^d$ be open and convex, and let $u:\Omega\to\R$ be convex
with Aleksandrov Monge--Ampère measure $M_u=0$ in $\Omega$.  Let
$x_0\in\Omega$ and let $p\in\partial u(x_0)$.  Then the contact set
\[
    C(x_0,p):=
    \{x\in\Omega:\ u(x)=u(x_0)+p\cdot(x-x_0)\}
\]
is not reduced to $\{x_0\}$ in any compactly contained convex neighborhood of
$x_0$.  Consequently the graph of $u$ contains a nontrivial affine segment
through $(x_0,u(x_0))$.
\end{proposition}

\begin{proof}
Subtract the supporting affine function and assume
\[
    u(x_0)=0,\qquad u\geq0,\qquad p=0.
\]
Suppose, toward a contradiction, that for some bounded convex
$U\Subset\Omega$ containing $x_0$ the contact set $\{u=0\}\cap U$ is exactly
$\{x_0\}$.  Then $u>0$ on $\partial U$.  Choose
\[
    0<h<\min_{\partial U}u.
\]
The section
\[
    S_h:=\{x\in U:\ u(x)<h\}
\]
is a nonempty bounded convex set with compact closure in $U$, and the convex
function
\[
    w:=u-h
\]
satisfies $w<0$ in $S_h$ and $w=0$ on $\partial S_h$.  Aleksandrov's maximum
principle gives, for every $x\in S_h$,
\[
    |w(x)|^d
    \leq C_d\,{\rm diam}(S_h)^{d-1}\,
        {\rm dist}(x,\partial S_h)\,M_u(S_h).
\]
But $M_u(S_h)=0$, hence $w\equiv0$ in $S_h$, contradicting
$w(x_0)=-h<0$.  Therefore every such neighborhood contains a second contact
point $x_1\neq x_0$.  Since $C(x_0,p)$ is convex, the segment
$[x_0,x_1]$ lies in $C(x_0,p)$, and the corresponding segment of the graph is
affine.
\end{proof}

\begin{remark}[Other rigidity thresholds]
There may also be a Sobolev or fractional $C^1$ route.  The role of this
option is different from the two alternatives above: it does not give a
canonical curvature measure below $C^1$, but it places the weak homogeneous
Monge--Ampère equation in a known rigidity regime.  In the two-dimensional weak
Monge--Ampère literature, developability for the degenerate equation is rigid
above certain thresholds, while convex-integration flexibility appears below
them.

For example, Jerrard's theorem is a rough rigidity result.  Lewicka--Pakzad's
convex-integration work exhibits the opposite, flexible side of the theory
\cite{LP17}.  The later De Lellis--Pakzad theorem proves developability for
very weak homogeneous Monge--Ampère solutions in the little Hölder class
$c^{1,2/3}$, hence for $C^{1,\alpha}$ with $\alpha>2/3$ \cite{DLP24}.  In the
Sobolev direction, a natural hypothesis to test is
\[
    \psi\in W^{2,2}_{\rm loc}.
\]
Then the equation $\det D^2\psi=0$ can be interpreted weakly, and the same
section-law argument above should force that weak equation.  If one then has a
developability theorem in this Sobolev class, the conclusion is again that the
graph is locally ruled by affine segments.  A higher-dimensional version would
need enough Sobolev control to make the nullity distribution of $D^2\psi$
integrable; this is precisely the step that fails in the flexible
convex-integration regime.  Thus this third route is plausible, but it should
be treated as a rigidity-threshold alternative rather than as the primary
low-regularity convex hypothesis.
\end{remark}

\subsection{Horizontal rulings and the final volume contradiction}
The nonnegativity normalization gives the last geometric restriction on a
complete ruling.  Suppose that the graph contains a complete affine line
\[
    \ell(s)=(x_0+s v,\psi(x_0)+s\alpha),\qquad s\in\R,
\]
with $v\in\R^d$ and $\alpha\in\R$.  Since $\ell$ lies in the graph, one has
\[
    \psi(x_0+s v)=\psi(x_0)+s\alpha
    \qquad\text{for every }s\in\R.
\]
If $\alpha\neq0$, then the right-hand side is negative for one of the two
limits $s\to+\infty$ or $s\to-\infty$, contradicting $\psi\geq0$.  Hence
\[
    \alpha=0.
\]
Thus every complete ruling is horizontal, i.e. parallel to the base
$\R^d\times\{0\}$, and $\psi$ is constant along the corresponding complete
line in the base.

This observation closes the argument as soon as the developability theorem
gives a cylindrical conclusion, namely a nontrivial linear subspace
$L\subset\R^d$ such that
\[
    \psi(x+z)=\psi(x)
    \qquad\text{for every }x\in\R^d,\ z\in L.
\]
Indeed, write $\R^d=L\oplus L^\perp$.  For $z\in L$ and $\eta\in L^\perp$,
the invariance gives
\[
    \Psi(2x,z+\eta)
    =\psi(x+\eta)+\psi(x-\eta)-2\psi(x).
\]
Consequently, for each $t>0$,
\[
    \{z+\eta:\ z\in L,\ \eta\in L^\perp,\
      \psi(x+\eta)+\psi(x-\eta)-2\psi(x)<t\}
    \subset \{y:\Psi(2x,y)<t\}.
\]
The transverse set in $L^\perp$ contains a neighborhood of $0$ by continuity
of $\psi$, while $L$ has infinite Lebesgue measure.  Hence
\[
    |\{y:\Psi(2x,y)<t\}|=+\infty
    \qquad\text{for every }t>0,
\]
contradicting the exact section law \eqref{eq:section-law}.

Thus the final higher-dimensional contradiction is reduced to the following
purely geometric point: the complete horizontal rulings obtained from zero
curvature should be parallel, or more generally should contain a common
nontrivial translation subspace $L$ along which $\psi$ is invariant.  This is
exactly what the Hartman--Nirenberg cylindrical conclusion gives in the smooth
complete case.  In the low-regularity alternatives above, this is the global
developability/cylinder statement that must be supplied.  Mere existence of one
complete horizontal line through each graph point gives only a zero-measure
line inside each centered section; parallelism is what upgrades this to an
infinite-volume product.

\begin{proposition}[Cylindrical obstruction]\label{prop:cylindrical-obstruction}
Let $d\geq3$, let $\psi:\R^d\to[0,\infty)$ be finite and convex, and assume
\eqref{eq:section-law}.  Suppose that there is a nontrivial linear subspace
$L\subset\R^d$ such that
\begin{equation}\label{eq:common-invariance}
    \psi(x+z)=\psi(x)
    \qquad\text{for every }x\in\R^d,\ z\in L.
\end{equation}
Then \eqref{eq:section-law} is impossible.
\end{proposition}

\begin{proof}
Fix $x\in\R^d$ and $t>0$.  Choose an orthogonal splitting
$\R^d=L\oplus L^\perp$.  For $z\in L$ and $\eta\in L^\perp$,
\eqref{eq:common-invariance} gives
\[
    \Psi(2x,z+\eta)
    =
    \psi(x+z+\eta)+\psi(x-z-\eta)-2\psi(x)
    =
    \psi(x+\eta)+\psi(x-\eta)-2\psi(x).
\]
The function
\[
    \eta\mapsto \psi(x+\eta)+\psi(x-\eta)-2\psi(x)
\]
is continuous on $L^\perp$ and equals $0$ at $\eta=0$.  Hence there is
$\rho>0$ such that it is $<t$ whenever $\eta\in L^\perp$ and $|\eta|<\rho$.
Thus, for every $R>0$,
\[
    \{z+\eta:\ z\in L,\ |z|<R,\ \eta\in L^\perp,\ |\eta|<\rho\}
    \subset \{y:\Psi(2x,y)<t\}.
\]
By the product decomposition of Lebesgue measure along
$L\oplus L^\perp$, the left-hand side has measure
\[
    \omega_{\dim L}R^{\dim L}\,\omega_{d-\dim L}\rho^{d-\dim L}.
\]
Letting $R\to\infty$ shows that
\[
    |\{y:\Psi(2x,y)<t\}|=+\infty,
\]
contradicting \eqref{eq:section-law}.
\end{proof}

\begin{corollary}[Conclusion under the Pogorelov alternative]
\label{cor:pogorelov-alternative}
Let $d\geq3$, let $\psi:\R^d\to[0,\infty)$ be finite and convex, and assume
\eqref{eq:section-law} and shrinking centered sections.  Assume moreover that
the graph of $\psi$ belongs to the Aleksandrov--Pogorelov bounded extrinsic
curvature class and that the corresponding global developability theorem gives
the cylindrical conclusion \eqref{eq:common-invariance} whenever the curvature
measure vanishes.  Then no such $\psi$ exists.
\end{corollary}

\begin{proof}
By the viscosity argument in the preceding subsection, the convolution identity
forces the Aleksandrov Monge--Ampère measure to vanish.  In the
Aleksandrov--Pogorelov class this is equivalent, through the spherical-image
parametrization of the graph, to vanishing extrinsic curvature measure on the
graph.  The assumed global Pogorelov developability theorem then gives a
nontrivial common translation subspace $L$ satisfying
\eqref{eq:common-invariance}.  Proposition~\ref{prop:cylindrical-obstruction}
contradicts \eqref{eq:section-law}.
\end{proof}

\begin{corollary}[Conclusion under the weak Monge--Ampère alternative]
\label{cor:ma-alternative}
Let $d\geq3$, let $\psi:\R^d\to[0,\infty)$ be finite and convex, and assume
\eqref{eq:section-law} and shrinking centered sections.  Assume moreover that
$\psi$ belongs to a weak Monge--Ampère class in which:
\begin{enumerate}
    \item the generalized determinant agrees with the Aleksandrov
    Monge--Ampère measure for convex functions;
    \item zero generalized determinant implies the global cylindrical
    conclusion \eqref{eq:common-invariance}.
\end{enumerate}
Then no such $\psi$ exists.
\end{corollary}

\begin{proof}
The argument above gives
\[
    {\rm Det}\,D^2\psi=0
\]
as a measure/current.  Proposition~\ref{prop:local-rulings-ma} gives the local
ruled conclusion at every point.  The second assumption in the present
corollary is the global rigidity input that promotes these local contact
segments to a common cylindrical translation subspace $L\neq\{0\}$.
Proposition~\ref{prop:cylindrical-obstruction} then gives the contradiction.
\end{proof}

For the present convolution problem, the first option seems best aligned with
convexity.  The regularity condition should be stated as membership in a
curvature-measure class for the graph, rather than as
$\det D^2\psi=0$ a.e.; the latter misses singular Monge--Ampère mass, exactly
as in the example $|x_1|+\frac12|x|^2$ discussed earlier.  If the forced
vanishing of this measure implies a global cylindrical ruling, the horizontal
ruling argument above rules out \eqref{eq:section-law} for $d\geq3$.  The
minimum regularity condition I would try first is therefore:
\[
\begin{gathered}
    \psi \text{ finite convex on }\R^d,\\
    \text{the graph of }\psi\text{ lies in the Aleksandrov--Pogorelov}\\
    \text{bounded extrinsic curvature class}.
\end{gathered}
\]
The convolution identity should then force the corresponding curvature measure
to be zero.  This is milder than $C^1$, keeps the singular curvature mass
visible, and is precisely tailored to the mechanism that should produce
rulings.

\begin{remark}
This is the point at which the two earlier drafts fit together.  The transcript
correctly identified that the implication
\[
    \det D^2\psi=\text{constant a.e.}
\]
does not exclude singular Alexandrov--Monge--Ampère mass.  The viscosity
comparison above is stronger: upper supports see the missing lower bound for
the determinant, and lower supports see the upper bound.  Thus no separate
uniqueness theorem for quadratic blowups is needed once the exact model
section law is available.
\end{remark}

\begin{remark}
For comparison, a positive definite quadratic graph in dimension $d$ has
centered convolution proportional to $\tau^{d/2-1}$.  Thus literal constancy
on the normalized support is compatible with a nondegenerate quadratic model
only when $d=2$.
\end{remark}

\section*{Acknowledgements}
DOS  is partially partially funded by Fundação para a Ciência e Tecnologia (FCT), Portugal, through  grant No.\@ UID/4459/2025 and  by the Deutsche Forschungsgemeinschaft 
 (DFG, German Research Foundation) under Germany's Excellence Strategy -- EXC-2047/1 -- 390685813.
JPR ...

\begin{thebibliography}{03}

\bibitem{CL03} L. Caffarelli, Y. Li,
\newblock {\it An extension to a theorem of Jörgens, Calabi, and Pogorelov.}
\newblock Comm. Pure Appl. Math. {\bf 56} (2003), no.~5, 549--583.

\bibitem{Caf90} L. Caffarelli,
\newblock {\it A localization property of viscosity solutions to the Monge--Ampère equation and their strict convexity.}
\newblock Ann. of Math. (2) {\bf 131} (1990), no.~1, 129--134.

\bibitem{Dud77} R. M. Dudley, 
\newblock {\it On second derivatives of convex functions.}
\newblock Math. Scand. {\bf 41} (1977), no.~1, 159--174.

\bibitem{DLP24} C. De Lellis, M. R. Pakzad,
\newblock {\it The geometry of $C^{1,\alpha}$ flat isometric immersions.}
\newblock Proc. Roy. Soc. Edinburgh Sect. A (2024), 1--39.

\bibitem{Fo99} G. B. Folland, 
\newblock Real Analysis.
Modern Techniques and Their Applications. Second edition
\newblock Pure Appl. Math. (N. Y.)
Wiley-Intersci. Publ.
John Wiley \& Sons, Inc., New York, 1999.


\bibitem{Gu16} C. Gutiérrez,
\newblock The Monge-Ampère equation.
Second edition. 
\newblock Progr. Nonlinear Differential Equations Appl., {\bf 89}
Birkhäuser/Springer, [Cham], 2016. 

\bibitem{HN59} P. Hartman, L. Nirenberg,
\newblock {\it On spherical image maps whose Jacobians do not change sign.}
\newblock Amer. J. Math. {\bf 81} (1959), 901--920.

\bibitem{Je10} R. L. Jerrard,
\newblock {\it Some rigidity results related to Monge--Ampère functions.}
\newblock Canad. J. Math. {\bf 62} (2010), no.~2, 320--354.

\bibitem{LP17} M. Lewicka, M. R. Pakzad,
\newblock {\it Convex integration for the Monge--Ampère equation in two dimensions.}
\newblock Anal. PDE {\bf 10} (2017), no.~3, 695--727.

\bibitem{OSQ18} D. Oliveira e Silva, R. Quilodrán, 
\newblock {\it On extremizers for Strichartz estimates for higher order Schrödinger equations.}
\newblock Trans. Amer. Math. Soc. {\bf 370} (2018), no.~10, 6871--6907.

\bibitem{OSQ20} D. Oliveira e Silva, R. Quilodrán, 
\newblock {\it A comparison principle for convolution measures with applications..}
\newblock Math. Proc. Cambridge Philos. Soc.   {\bf 169} (2020), no.~2, 307--322.

\bibitem{Po73} A. V. Pogorelov,
\newblock {\it Extrinsic Geometry of Convex Surfaces.}
\newblock Translations of Mathematical Monographs, vol.~35,
American Mathematical Society, Providence, RI, 1973.




\bibitem{St93} E. M. Stein,
\newblock Harmonic Analysis: Real-Variable Methods, Orthogonality, and Oscillatory Integrals.
\newblock Princeton Univ. Press, Princeton, NJ, 1993.

\bibitem{Us99} V. Ushakov,
\newblock {\it Developable surfaces in Euclidean space.}
\newblock J. Aust. Math. Soc. {\bf 66} (1999), no.~3, 388--402.
\end{thebibliography}


\end{document}  
